% GENERATED FILE. Edit canonical lessons, then run npm run generate:books.
% canonical-source-hash: fnv1a64:b8fd6e12f0bf7e7a
% canonical-lessons: HI-C303-hilna, HI-C303-jhatka, HI-C303-darar, HI-C303-dhasna, HI-C303-dabna, HI-R303-review-296-303

\chapter{When the Ground Shook}
\label{ch:hi-when-the-ground-shook}
\hlchaptermodality{303}

\begin{chapteropening}
\textbf{By the end of this chapter:} \emph{I can tell what happened in an earthquake.}

Complete the last lesson of chapter 303: I can tell what happened in an earthquake.
\end{chapteropening}

\section[\texorpdfstring{\hi{हिलना} (hilnā)}{हिलना (hilnā)}]{\hi{हिलना} (hilnā) — to shake (by itself); to move}

\label{lesson:HI-C303-hilna}



\begin{quote}
Before the new one: say the Hindi for careful, then the Hindi for to be averted.
\end{quote}

\subsection*{You'll want to know: \hi{हिलना}}
\textbf{\hi{हिलना}} — \emph{hilnā} — \textquotedblleft{}to shake (by itself); to move\textquotedblright{}.

\begin{quote}
\hi{रात} \hi{को} \hi{अचानक} \hi{पूरा} \hi{घर} \hi{हिला।}
\end{quote}

\begin{quote}
\emph{rāt ko acānak pūrā ghar hilā.}
\end{quote}

\begin{quote}
Suddenly, in the night, the whole house shook.
\end{quote}

Past (it …): \hi{वह} \hi{पेड़} \hi{हिला}, \emph{vah peṛ hilā}.

\begin{cousinweb}
\textbf{\hi{हिलना}} (\emph{hilnā}) is inherited from \emph{*hillati}, \textquotedblleft{}shakes\textquotedblright{}, a form CDIAL reconstructs (CDIAL 14120). It is the intransitive partner of \textbf{\hi{हिलाना}}: you \textbf{\hi{हिलाते} \hi{हैं}} a branch, and the branch \textbf{\hi{हिलती} \hi{है}}.
\end{cousinweb}

\begin{grammarlens}[title={What you've built}]
One more word for telling what happened.
\end{grammarlens}

\subsection*{Your turn}
\begin{itemize}
\raggedright
  \item \emph{Say it:} \emph{hilnā}
  \item \emph{Say it:} \emph{hilnā}, once more
  \item \emph{Say it:} \emph{rāt ko acānak pūrā ghar hilā.} — Suddenly, in the night, the whole house shook.
  \item \emph{Recall:} say the Hindi for careful, then the Hindi for to be averted, then say \emph{hilnā} again
\end{itemize}

\begin{tcolorbox}[breakable,colback=teal!4,colframe=teal!35!black,title={Before you move on}]
What does \hi{हिलना} mean? (\textquotedblleft{}To shake (by itself); to move\textquotedblright{}.) Say it once more.
\end{tcolorbox}
\section[\texorpdfstring{\hi{झटका} (jhaṭkā)}{झटका (jhaṭkā)}]{\hi{झटका} (jhaṭkā) — a jolt; a tremor}

\label{lesson:HI-C303-jhatka}



\begin{quote}
Before the new one: say the Hindi for to be averted, then the Hindi for to shake (by itself).
\end{quote}

\subsection*{You'll want to know: \hi{झटका}}
\textbf{\hi{झटका}} — \emph{jhaṭkā} — \textquotedblleft{}a jolt; a tremor\textquotedblright{}.

\begin{quote}
\hi{भूकंप} \hi{का} \hi{पहला} \hi{झटका} \hi{सुबह} \hi{चार} \hi{बजे} \hi{आया।}
\end{quote}

\begin{quote}
\emph{bhūkamp kā pahlā jhaṭkā subah cār baje āyā.}
\end{quote}

\begin{quote}
The first tremor of the earthquake came at four in the morning.
\end{quote}

Masculine.

\begin{cousinweb}
is inherited from \emph{*jhaṭ}, a form CDIAL reconstructs for quick, sudden movements (CDIAL 5327). An earthquake comes in \textbf{\hi{झटके}}; bad news can be a \textbf{\hi{झटका}} too.
\end{cousinweb}

\begin{grammarlens}[title={What you've built}]
One more word for telling what happened.
\end{grammarlens}

\subsection*{Your turn}
\begin{itemize}
\raggedright
  \item \emph{Say it:} \emph{jhaṭkā}
  \item \emph{Say it:} \emph{jhaṭkā}, once more
  \item \emph{Say it:} \emph{bhūkamp kā pahlā jhaṭkā subah cār baje āyā.} — The first tremor of the earthquake came at four in the morning.
  \item \emph{Recall:} say the Hindi for to be averted, then the Hindi for to shake (by itself), then say \emph{jhaṭkā} again
\end{itemize}

\begin{tcolorbox}[breakable,colback=teal!4,colframe=teal!35!black,title={Before you move on}]
What does \hi{झटका} mean? (\textquotedblleft{}A jolt; a tremor\textquotedblright{}.) Say it once more.
\end{tcolorbox}
\section[\texorpdfstring{\hi{दरार} (darār)}{दरार (darār)}]{\hi{दरार} (darār) — a crack}

\label{lesson:HI-C303-darar}



\begin{quote}
Before the new one: say the Hindi for to shake (by itself), then the Hindi for a jolt.
\end{quote}

\subsection*{You'll want to know: \hi{दरार}}
\textbf{\hi{दरार}} — \emph{darār} — \textquotedblleft{}a crack\textquotedblright{}.

\begin{quote}
\hi{झटके} \hi{से} \hi{दीवार} \hi{में} \hi{दरार} \hi{आ} \hi{गई।}
\end{quote}

\begin{quote}
\emph{jhaṭke se dīvār meṁ darār ā gaī.}
\end{quote}

\begin{quote}
The jolt left a crack in the wall.
\end{quote}

Feminine.

\begin{cousinweb}
\textbf{\hi{दरार}} (\emph{darār}) is listed by CDIAL under \emph{*darākāra}, \textquotedblleft{}shaped like a crack\textquotedblright{}, a reconstructed form (CDIAL 6193). A crack \textbf{\hi{आती} \hi{है}} or \textbf{\hi{पड़ती} \hi{है}} in a wall.
\end{cousinweb}

\begin{grammarlens}[title={What you've built}]
One more word for telling what happened.
\end{grammarlens}

\subsection*{Your turn}
\begin{itemize}
\raggedright
  \item \emph{Say it:} \emph{darār}
  \item \emph{Say it:} \emph{darār}, once more
  \item \emph{Say it:} \emph{jhaṭke se dīvār meṁ darār ā gaī.} — The jolt left a crack in the wall.
  \item \emph{Recall:} say the Hindi for to shake (by itself), then the Hindi for a jolt, then say \emph{darār} again
\end{itemize}

\begin{tcolorbox}[breakable,colback=teal!4,colframe=teal!35!black,title={Before you move on}]
What does \hi{दरार} mean? (\textquotedblleft{}A crack\textquotedblright{}.) Say it once more.
\end{tcolorbox}
\section[\texorpdfstring{\hi{धँसना} (dhãsnā)}{धँसना (dhãsnā)}]{\hi{धँसना} (dhãsnā) — to cave in; to sink in}

\label{lesson:HI-C303-dhasna}



\begin{quote}
Before the new one: say the Hindi for a jolt, then the Hindi for a crack.
\end{quote}

\subsection*{You'll want to know: \hi{धँसना}}
\textbf{\hi{धँसना}} — \emph{dhãsnā} — \textquotedblleft{}to cave in; to sink in\textquotedblright{}.

\begin{quote}
\hi{भूकंप} \hi{में} \hi{सड़क} \hi{नीचे} \hi{धँस} \hi{गई।}
\end{quote}

\begin{quote}
\emph{bhūkamp meṁ saṛak nīce dhãs gaī.}
\end{quote}

\begin{quote}
In the earthquake the road caved in.
\end{quote}

Past (it …): \hi{वह} \hi{सड़क} \hi{धँस} \hi{गई}, \emph{vah saṛak dhãs gaī}.

\begin{cousinweb}
\textbf{\hi{धँसना}} (\emph{dhãsnā}) goes back to Sanskrit \emph{dhvaṁsati}, \textquotedblleft{}falls to pieces\textquotedblright{} (CDIAL 6896). Ground, a road or a foot in mud \textbf{\hi{धँसता} \hi{है}}: it goes down into what is under it.
\end{cousinweb}

\begin{grammarlens}[title={What you've built}]
One more word for telling what happened.
\end{grammarlens}

\subsection*{Your turn}
\begin{itemize}
\raggedright
  \item \emph{Say it:} \emph{dhãsnā}
  \item \emph{Say it:} \emph{dhãsnā}, once more
  \item \emph{Say it:} \emph{bhūkamp meṁ saṛak nīce dhãs gaī.} — In the earthquake the road caved in.
  \item \emph{Recall:} say the Hindi for a jolt, then the Hindi for a crack, then say \emph{dhãsnā} again
\end{itemize}

\begin{tcolorbox}[breakable,colback=teal!4,colframe=teal!35!black,title={Before you move on}]
What does \hi{धँसना} mean? (\textquotedblleft{}To cave in; to sink in\textquotedblright{}.) Say it once more.
\end{tcolorbox}
\section[\texorpdfstring{\hi{दबना} (dabnā)}{दबना (dabnā)}]{\hi{दबना} (dabnā) — to be pressed down; to be trapped under}

\label{lesson:HI-C303-dabna}



\begin{quote}
Before the new one: say the Hindi for a crack, then the Hindi for to cave in.
\end{quote}

\subsection*{You'll want to know: \hi{दबना}}
\textbf{\hi{दबना}} — \emph{dabnā} — \textquotedblleft{}to be pressed down; to be trapped under\textquotedblright{}.

\begin{quote}
\hi{दीवार} \hi{गिरी} \hi{और} \hi{दो} \hi{बकरियाँ} \hi{उसके} \hi{नीचे} \hi{दब} \hi{गईं।}
\end{quote}

\begin{quote}
\emph{dīvār girī aur do bakriyā̃ uske nīce dab gaīṁ.}
\end{quote}

\begin{quote}
The wall fell and two goats were trapped under it.
\end{quote}

Past (he …): \hi{वह} \hi{आदमी} \hi{दब} \hi{गया}, \emph{vah ādmī dab gayā}.

\begin{cousinweb}
\textbf{\hi{दबना}} (\emph{dabnā}) is inherited from \emph{*dabb}, \textquotedblleft{}press\textquotedblright{}, a form CDIAL reconstructs (CDIAL 6173). It is the intransitive partner of \textbf{\hi{दबाना}}, to press: what is pressed down \textbf{\hi{दबता} \hi{है}}.
\end{cousinweb}

\begin{grammarlens}[title={What you've built}]
One more word for telling what happened.
\end{grammarlens}

\subsection*{Your turn}
\begin{itemize}
\raggedright
  \item \emph{Say it:} \emph{dabnā}
  \item \emph{Say it:} \emph{dabnā}, once more
  \item \emph{Say it:} \emph{dīvār girī aur do bakriyā̃ uske nīce dab gaīṁ.} — The wall fell and two goats were trapped under it.
  \item \emph{Recall:} say the Hindi for a crack, then the Hindi for to cave in, then say \emph{dabnā} again
\end{itemize}

\begin{tcolorbox}[breakable,colback=teal!4,colframe=teal!35!black,title={Before you move on}]
What does \hi{दबना} mean? (\textquotedblleft{}To be pressed down; to be trapped under\textquotedblright{}.) Say it once more.
\end{tcolorbox}
\section[(dialogue)]{Looking back: from A Crash on the Road to When the Ground Shook}

\label{lesson:HI-R303-review-296-303}



\begin{quote}
Say the Hindi for speed and to collide.
\end{quote}

\subsection*{Your turn}
Say these aloud:

\begin{itemize}
\raggedright
  \item speed, to collide, a passenger, a wheel, to overturn
  \item to slip, a stumble, to stagger, to be grazed, covered in blood
  \item a mishap, injured, unconscious, to crush, to steady oneself
  \item an ember, to smoulder, a flame, to flare up, dry straw
  \item to blaze, to glow red-hot, a blaze, light, scorching heat
  \item a siren, to be flustered, to go out (of a fire or light), to melt, a blister
  \item a disaster, a warning, an announcement, careful, to be averted
  \item to shake (by itself), a jolt, a crack, to cave in, to be pressed down
\end{itemize}

\begin{tcolorbox}[breakable,colback=teal!4,colframe=teal!35!black,title={Before you move on}]
Speed? (\textbf{\hi{रफ़्तार}}, \emph{raftār}.) To collide? (\textbf{\hi{टकराना}}, \emph{ṭakrānā}.) A passenger? (\textbf{\hi{सवारी}}, \emph{savārī}.) A wheel? (\textbf{\hi{पहिया}}, \emph{pahiyā}.) To overturn? (\textbf{\hi{पलटना}}, \emph{palaṭnā}.) To slip? (\textbf{\hi{फिसलना}}, \emph{phisalnā}.) A stumble? (\textbf{\hi{ठोकर}}, \emph{ṭhokar}.) To stagger? (\textbf{\hi{लड़खड़ाना}}, \emph{laṛkhaṛānā}.) To be grazed? (\textbf{\hi{छिलना}}, \emph{chilnā}.) Covered in blood? (\textbf{\hi{लहूलुहान}}, \emph{lahūluhān}.) A mishap? (\textbf{\hi{हादसा}}, \emph{hādsā}.) Injured? (\textbf{\hi{घायल}}, \emph{ghāyal}.) Unconscious? (\textbf{\hi{बेहोश}}, \emph{behoś}.) To crush? (\textbf{\hi{कुचलना}}, \emph{kucalnā}.) To steady oneself? (\textbf{\hi{सँभलना}}, \emph{sãbhalnā}.) An ember? (\textbf{\hi{अंगारा}}, \emph{aṅgārā}.) To smoulder? (\textbf{\hi{सुलगना}}, \emph{sulagnā}.) A flame? (\textbf{\hi{शोला}}, \emph{śolā}.) To flare up? (\textbf{\hi{भड़कना}}, \emph{bhaṛaknā}.) Dry straw? (\textbf{\hi{फूस}}, \emph{phūs}.) To blaze? (\textbf{\hi{धधकना}}, \emph{dhadhaknā}.) To glow red-hot? (\textbf{\hi{दहकना}}, \emph{dahaknā}.) A blaze? (\textbf{\hi{ज्वाला}}, \emph{jvālā}.) Light? (\textbf{\hi{उजाला}}, \emph{ujālā}.) Scorching heat? (\textbf{\hi{तपिश}}, \emph{tapiś}.) A siren? (\textbf{\hi{सायरन}}, \emph{sāyran}.) To be flustered? (\textbf{\hi{हड़बड़ाना}}, \emph{haṛbaṛānā}.) To go out (of a fire or light)? (\textbf{\hi{बुझना}}, \emph{bujhnā}.) To melt? (\textbf{\hi{पिघलना}}, \emph{pighalnā}.) A blister? (\textbf{\hi{फफोला}}, \emph{phapholā}.) A disaster? (\textbf{\hi{आपदा}}, \emph{āpdā}.) A warning? (\textbf{\hi{चेतावनी}}, \emph{cetāvnī}.) An announcement? (\textbf{\hi{ऐलान}}, \emph{ailān}.) Careful? (\textbf{\hi{सावधान}}, \emph{sāvdhān}.) To be averted? (\textbf{\hi{टलना}}, \emph{ṭalnā}.) To shake (by itself)? (\textbf{\hi{हिलना}}, \emph{hilnā}.) A jolt? (\textbf{\hi{झटका}}, \emph{jhaṭkā}.) A crack? (\textbf{\hi{दरार}}, \emph{darār}.) To cave in? (\textbf{\hi{धँसना}}, \emph{dhãsnā}.) To be pressed down? (\textbf{\hi{दबना}}, \emph{dabnā}.)
\end{tcolorbox}
